% !TeX root = ../main.tex
\begin{frame}{03 创新思路——单胞设计}
  \begin{SlideCanvas}
    \BodyAt{55}{84}{850}{本课题组设计反手性单胞，通过周期阵列与双向拉伸计算检验目标应力、切线刚度及泊松比的匹配。}
    \PicAt{55}{145}{592}{294}{assets/group_singlecell.pdf}
    \PointAt{688}{161}{214}{反手性构型}
    \BodyAt{688}{207}{214}{$\nu^\ast=-0.8$\\[12pt]
      $\varepsilon_{\max}=30\%$}
    \PointAt{688}{316}{214}{超弹性表征}
    \DetailAt{688}{355}{214}{等效响应来自结构变形，基体材料无需具有负泊松比。}
    \DetailAt{55}{445}{590}{图(e)—(g)分别比较应力、切线刚度与泊松比的目标拟合。}
    \Sources{本课题组成果：\RefGroup，图4。几何、坐标与曲线保留原文标注。}
  \end{SlideCanvas}
  \BICConclusionBand{设计评价覆盖\BICBandEmph{有限变形区间}，并检查目标曲线的拟合误差}
  \note{图4(a)为设计单胞，(b)为5×5阵列，(c)(d)为两个拉伸方向的构型，(e)(f)(g)为应力、切线刚度和泊松比曲线。论文示例的单胞最大拉伸应变为30\%。论文采用割线应变比描述泊松比，已发表算例主要展示负值。用户说明方法也能设计零值，后续支架目标将其纳入候选，并在对应周向工况验证。图4不代表真实支架部署试验。}
\end{frame}
